\documentclass[10pt,a4paper]{amsart}
\usepackage[margin=19mm]{geometry}
\usepackage{amsmath,amssymb,mathtools}
\usepackage[scr=rsfs,cal=euler]{mathalfa}
\usepackage{xfrac}
\usepackage[unicode,hidelinks,bookmarks=false]{hyperref}
\newcommand{\ie}{i.e.\ }
\newcommand{\cf}{cf.\ }
\newcommand{\resp}{respectively\ }
\newcommand{\Subsection}{\subsection}
\providecommand{\nobreakdash}{\nobreak\hbox{-}}
\setlength{\emergencystretch}{2em}
\multlinegap0pt
\newcommand{\C}{\mathbb{C}}
\newcommand{\dd}{\,\text{d}}
\newcommand{\pr}{\partial}
\newcommand{\psiz}{\psi_{z_0}}
\theoremstyle{plain}
\newtheorem{thm}{Theorem}[section]
\newtheorem{prop}[thm]{Proposition}
\newtheorem{lem}[thm]{Lemma}
\newtheorem{cor}[thm]{Corollary}
\theoremstyle{definition}
\newtheorem{dfn}[thm]{Definition}
\theoremstyle{remark}
\newtheorem{rem}[thm]{Remark}
\title[Uniqueness for the inverse Schr\"odinger problem in the plane with $L^p$ potentials, $p>1$]{Uniqueness for the inverse Schr\"odinger problem in the plane with $L^p$ potentials, $p>1$}
\author{C\u{a}t\u{a}lin I. C\^arstea, Jenn-Nan Wang}
\date{Source: arXiv:2609.28046v1 (CC BY 4.0)}
\begin{document}
\maketitle

\noindent\emph{Source and licence.} The delivered work is the article shown in the title above, version arXiv:2609.28046v1 (CC BY 4.0), distributed under the Creative Commons Attribution 4.0 International licence, as recorded in the arXiv licence metadata for this exact version and shown on the abstract page saved with this item. The full licence notice, which carries the licence URL and the address of that abstract page, is the bundled file \texttt{licenses/arxiv-2609.28046-CC-BY-4.0.txt}.
\par\smallskip

\noindent\textit{Editorial scope.} Counted unit: the article's Lemma (source label lem-inf) (source0.tex lines 529--535, source label \texttt{lem-inf}): ``Let and let be compact. There is a constant such that, uniformly for and , -1 f L (X) C - 1 2 (2 ) f W 1,b (X) , f W 1,b 0(X), where is viewed as extended by zero outside .''. It is delivered together with the article's complete original proof (source0.tex lines 536--637, one proof environment adjacent to the statement, 102 lines), copied line for line from the article's own text. Required context retained uncounted: none beyond the counted statement and proof. Editorial formatting only: an AMS \texttt{amsart} preamble that restates the article's own macro definitions and its own theorem declarations, and the theorem environments the retained lines use; the article's own class file is neither shipped nor input; the retained environments are renumbered sequentially in order of appearance, whereas the article prints them with its own numbering; the article's equation numbering is not reproduced and every cross-reference inside the retained text resolves to the wrapper's numbering. No citation occurs inside any retained line, so the wrapper carries no bibliography; All retained bodies are exact copies of the bundled \texttt{sources/211548/source0.tex}, so no omission receipt applies. No asset-inclusion command occurs inside any retained span and the package ships no image asset.


\section{The counted unit: Lemma (source label lem-inf) and its complete original proof}

\begin{lem}\label{lem-inf}
Let $2<b<\infty$ and let $Z\Subset\C$ be compact.  There is a constant $C=C_{X,Z,b}>0$ such that, uniformly for $z_0\in Z$ and $\tau\geq2$,
\begin{equation}
\|\pr_\psi^{-1}f\|_{L^\infty(X)}\leq  C\tau^{-\frac{1}{2}}\log(2+\tau) \|f\|_{W^{1,b}(X)},\quad \forall f\in W^{1,b}_0(X),
\end{equation}
where $f$ is viewed as extended by zero outside $X$.
\end{lem}

\begin{proof}
We regard $f$ as extended by zero outside $X$ and put
\begin{equation}
        R_0=1+\sup\{|z-z_0|:z\in\overline{X},\ z_0\in Z\}<\infty.
\end{equation}
Choose $\chi\in C_0^\infty(\C)$ with $0\leq \chi\leq1$, equal to one on $B(0,1)$ and to zero outside $B(0,2)$.  For $\delta>0$, let $\chi_\delta(z)=\chi(\delta^{-1}(z-z_0))$.

We split $\pr_\psi^{-1}f$ into
\begin{equation}
\pr_\psi^{-1} f=\pr_\psi^{-1}(\chi_\delta f)+\pr_\psi^{-1}\left((1-\chi_\delta)f\right).
\end{equation}
The near-center term has the kernel representation
\begin{equation}
\pr_\psi^{-1}(\chi_\delta f)(z)=\frac{1}{\pi}\int \frac{e^{i\tau\psiz(z')}}{\bar z-\bar z'}\chi(\delta^{-1}(z'-z_0))f(z')\dd^2 z',
\end{equation}
which is bounded by
\begin{equation}
\|\pr_\psi^{-1}(\chi_\delta f)\|_{L^\infty(X)}\leq C\|f\|_{L^\infty(X)}\int_0^{2\delta}\dd\rho=C\delta \|f\|_{W^{1,b}(X)}.
\end{equation}

For the term away from the center, integration by parts in the oscillation gives
\begin{multline}
\pr_\psi^{-1}\left((1-\chi_\delta)f\right)
=\frac{1}{i\tau}\pr^{-1}\left[(\pr e^{i\tau\psiz}) \frac{1-\chi_\delta}{\pr\psiz} f  \right]\\[5pt]
=\frac{1}{i\tau}e^{i\tau\psiz}\frac{1-\chi_\delta}{\pr\psiz} f
-\frac{1}{i\tau}\pr_\psi^{-1}\left[\pr\left(\frac{1-\chi_\delta}{\pr\psiz} f\right)\right].
\end{multline}
Since $\partial\psi_{z_0}=z-z_0$ and $\partial^2\psi_{z_0}=1$, the derivative in the last term is
\begin{equation}\label{product-rule-oscillatory-cauchy}
\partial\left(\frac{1-\chi_\delta}{\partial\psi_{z_0}}f\right)
=\frac{1-\chi_\delta}{\partial\psi_{z_0}}\partial f
-\frac{\partial\chi_\delta}{\partial\psi_{z_0}}f
-\frac{1-\chi_\delta}{(\partial\psi_{z_0})^2}f .
\end{equation}
We estimate the three terms in \eqref{product-rule-oscillatory-cauchy} separately.  The cutoff distance gives
\begin{equation}
\left\Vert \frac{1-\chi_\delta}{\pr\psiz} f\right\Vert_{L^\infty(X)}\leq\frac{\|f\|_{L^\infty(X)}}{\delta},
\end{equation}
and the local Cauchy estimate similarly gives
\begin{equation}
\left\Vert\pr_\psi^{-1}\left[\frac{1-\chi_\delta}{\pr\psiz} \pr f  \right]\right\Vert_{L^\infty(X)}\leq C\frac{\|\pr f\|_{L^b(X)}}{\delta}.
\end{equation}
The derivative of the cutoff is supported on the annulus $\delta\leq |z'-z_0|\leq2\delta$, where
\begin{multline}
\left|\pr_\psi^{-1}\left[\frac{\pr \chi_\delta}{\pr\psiz} f  \right](z)\right| \\
\leq C\delta^{-1}\|f\|_{L^\infty(X)}
\int_{\delta\leq|z'-z_0|\leq2\delta}
\frac{\dd^2z'}{|z-z'|\,|z'-z_0|}\\[5pt]
= C\delta^{-1}\|f\|_{L^\infty(X)}I(z-z_0),
\end{multline}
where
\begin{align}
I(z)
&=\int_{\delta\leq|z'|\leq2\delta}\frac{1}{|z-z'|}\frac{1}{|z'|}\dd^2z' \notag\\
&=\int_{\delta\leq|z'|\leq2\delta}\mathbf{1}_{B_{z}(\delta)}(z')\frac{1}{|z-z'|}\frac{1}{|z'|}\dd^2z' \notag\\
&\quad +\int_{\delta\leq|z'|\leq2\delta}\mathbf{1}_{B_{z}(\delta)^c}(z')\frac{1}{|z-z'|}\frac{1}{|z'|}\dd^2z'
=I_1(z)+I_2(z).
\end{align}
The portion of this annulus within distance $\delta$ of $z$ satisfies
\begin{equation}
I_1(z)\leq\int_{|z'-z|\leq \delta}\frac{1}{\delta|z'-z|}\dd^2 z'\leq C
\end{equation}
whereas its complement satisfies
\begin{equation}
I_2(z)\leq \int_{|z'|\leq2\delta}\frac{1}{\delta|z'|}\dd^2 z'\leq C.
\end{equation}
Therefore
\begin{equation}
\left|\pr_\psi^{-1}\left[\frac{\pr \chi_\delta}{\pr\psiz} f  \right](z)\right|\leq C\delta^{-1}\|f\|_{L^\infty(X)}.
\end{equation}


For the square-denominator term, the same decomposition gives
\begin{multline}
\left|\pr_\psi^{-1}\left[\frac{1- \chi_\delta}{(\pr\psiz)^2} f  \right](z)\right|\\[5pt]
\leq C\|f\|_{L^\infty(X)}\int_{R_0\geq|z'-z_0|\geq\delta}\frac{1}{|z-z'|}\frac{1}{|z'-z_0|^2}\dd^2z'\\[5pt]
=C\|f\|_{L^\infty(X)}J(z-z_0),
\end{multline}
where 
\begin{align}
J(z)
&=\int_{\delta\leq|z'|\leq R_0}\frac{1}{|z-z'|}\frac{1}{|z'|^{2}}\dd^2z' \notag\\
&=\int_{\delta\leq|z'|\leq R_0}\mathbf{1}_{B_{z}(\delta)}(z')\frac{1}{|z-z'|}\frac{1}{|z'|^{2}}\dd^2z' \notag\\
&\quad +\int_{\delta\leq|z'|\leq R_0}\mathbf{1}_{B_{z}(\delta)^c}(z')\frac{1}{|z-z'|}\frac{1}{|z'|^{2}}\dd^2z'
=J_1(z)+J_2(z).
\end{align}
Here
\begin{equation}
J_1(z)\leq\int_{|z'-z|\leq \delta}\frac{1}{\delta^2|z'-z|}\dd^2 z'\leq C\delta^{-1}
\end{equation}
and
\begin{equation}
J_2(z)\leq \int_{R_0\geq|z'|\geq\delta}\frac{1}{\delta|z'|^2}\dd^2 z'\leq C\delta^{-1}|\log\delta|.
\end{equation}


Combining these estimates,
\begin{equation}
\|\pr_\psi^{-1}f\|_{L^\infty(X)}\leq C\|f\|_{W^{1,b}(X)}\left(\delta+\tau^{-1}\delta^{-1} |\log\delta| \right).
\end{equation}
Choosing $\delta=\tau^{-1/2}$ gives the stated estimate.
\end{proof}

\end{document}
